\scanpage{17}
\uline{Cohomologie à supports compacts} \deleted{de} de $C,E,\widehat E\ldots$

Les \unsure{cohomologies} à support compact de $E,\widehat E$ se \unsure{déduit} simplement par dualité des \deleted{celles}\inserted{la coh. usuelle} de $E,\widehat E$ lorsque $X$ est lisse \unsure{\ldots\ (\ldots}. Celle de $\widehat C$ est identique à sa cohomologie usuelle (car $\widehat C$ est propre\ldots).

Reste à déterminer la coh. à support compact de $C$.

On a
\[
\deleted{H_c^*(C')\longleftarrow H_c^{\unsure{*}}(\unsure{\,})}
\]
\diagram{\begin{array}{ccccccccccc}
\cdots&\longrightarrow&H^i(\widehat C)&\text{---}&H^i(X)&\text{---}&H_c^{i+1}(C)&\text{---}&H^{i+1}(\widehat C)&\text{---}&H^{i+1}(X)\\
&&\wr&&&&&&\wr\\
&&H^{i-2}(X)(-1)&&&&&&H^{i-1}(X)(-1)\\
&&\text{si }i\neq0
\end{array}}

Comparant avec la suite exacte analogue pour déterminer \inserted{la coh. de} $E$ \inserted{\unsure{en calculant} $H_c^*(\widehat C)$ par $\widehat E\subset\widehat C$}, \unsure{on est amené à} prévoir
\[
H^i(E)\underset{\text{si }i\neq0}{\simeq}H_c^{i+1}(C)
\simeq H_{\unsure{a}}^{i+1}(C),
\qquad \deleted{\unsure{et}}
\]
\deleted{\unsure{on a en effet, pour l'inclusion} $E\subset C$,}
\[
\boxed{H_a^*(C)\simeq H_c^*(C)\quad ??}
\]
\[
\deleted{H_a^i(C)=H_c^{\unsure{i}}(C)=}
\]
\mspage{16}
